\documentclass[10pt]{article}
\usepackage[T1]{fontenc}
\usepackage[utf8]{inputenc}
\usepackage[french]{babel}
\usepackage{amsmath,amssymb,mathrsfs,mathpazo}
\usepackage{geometry}
\geometry{a4paper,margin=2cm}
\usepackage{graphicx,array,tabularx,tabularray,colortbl,multirow,multicol,booktabs,enumitem,pifont,ulem,fancybox,diagbox,cancel}
\usepackage{pstricks,pstricks-add,pst-plot,pst-tree,pst-node,pst-text,pst-func,pst-math,pst-bspline,pst-3dplot,pst-eucl,pst-solides3d}
\let\pscalculate\undefined
\usepackage{pst-bezier}
\usepackage{tikz,pgfplots,tkz-tab}
\usetikzlibrary{arrows,arrows.meta,patterns,calc,positioning,trees,decorations.markings,shapes.geometric,shapes.misc,shapes.arrows,intersections,decorations.pathmorphing,decorations.pathreplacing,backgrounds,shadows,babel}
\pgfplotsset{compat=1.18}
\usepgfplotslibrary{fillbetween}
\usepackage[np]{numprint}
\usepackage{listings,makecell,pas-tableur,twoopt}
\let\barre\undefined
\usepackage[pstricks]{tablvar}
\definecolor{darkgreen}{RGB}{0,100,0}
\providecommand{\red}{\color{red}}
\providecommand{\blue}{\color{blue}}
\providecommand{\green}{\color{green}}
\providecommand{\black}{\color{black}}
\providecommand{\white}{\color{white}}
\providecommand{\gray}{\color{gray}}
\providecommand{\cyan}{\color{cyan}}
\providecommand{\magenta}{\color{magenta}}
\providecommand{\yellow}{\color{yellow}}
\providecommand{\violet}{\color{violet}}
\providecommand{\orange}{\color{orange}}
\providecommand{\euro}{\texteuro}
\providecommand{\EUR}{\texteuro}
\providecommand{\ieme}{\textsuperscript{e}}
\providecommand{\degres}{\ensuremath{^\circ}}
\providecommand{\degree}{\ensuremath{^\circ}}
\providecommand{\vv}[1]{\overrightarrow{#1}}
\providecommand{\Coord}[1]{#1}
\providecommand{\decosix}{}
\providecommand{\starredbullet}{\ensuremath{\bullet}}
\pagestyle{empty}
\setlength{\parindent}{0pt}
\usepackage{adjustbox}\DeclareUnicodeCharacter{2236}{:}\DeclareUnicodeCharacter{2190}{\ensuremath{\leftarrow}}
\providecommand{\R}{\mathbb{R}}
\providecommand{\N}{\mathbb{N}}
\providecommand{\Z}{\mathbb{Z}}
\providecommand{\Q}{\mathbb{Q}}
\providecommand{\D}{\mathbb{D}}
\providecommand{\C}{\mathbb{C}}
\providecommand{\e}{{\mathrm{e}}}
\providecommand{\vect}[1]{\overrightarrow{#1}}
\providecommand{\vectt}[1]{\overrightarrow{\mathrm{#1}}}
\providecommand{\vvt}[1]{\overrightarrow{\mathrm{#1}}}
\providecommand{\np}[2][]{{#2~#1}}
\providecommand{\vv}[1]{{\overrightarrow{#1}}}
\renewcommand{\d}{{\mathrm{d}}}
\providecommand{\degres}{{^\circ}}
\providecommand{\degre}{{^\circ}}
\providecommand{\blue}{\color{blue}}
\providecommand{\red}{\color{red}}
\providecommand{\enskip}{\quad}
\providecommand{\llbracket}{\lbrack\!\lbrack}
\providecommand{\rrbracket}{\rbrack\!\rbrack}
\providecommand{\euro}{{\text{€}}}
\providecommand{\no}{{\text{n°}}}
\providecommand{\interval}[2]{{[#1;#2]}}
\providecommand{\intervFO}[2]{{[#1;#2[}}
\providecommand{\dotfill}{\ldots}

\providecommand{\ds}{\displaystyle}
\providecommand{\textding}[1]{#1}
\providecommand{\Oij}{\(\mathrm{O};\vec{\imath},\vec{\jmath}\)}
\providecommand{\Oijk}{\(\mathrm{O};\vec{\imath},\vec{\jmath},\vec{k}\)}
\providecommand{\Ouv}{\(\mathrm{O};\vec{u},\vec{v}\)}
\providecommand{\intervFF}[2]{\ensuremath{[#1;#2]}}
\renewcommand{\euro}{\ifmmode\text{\texteuro}\else\texteuro\fi}
\DeclareUnicodeCharacter{00B0}{\ifmmode^{\circ}\else\textdegree\fi}
\DeclareUnicodeCharacter{20AC}{\ifmmode\text{\texteuro}\else\texteuro\fi}
% Énoncés : Éducation nationale. Transcription et corrigés : APMEP, Denis Vergès et contributeurs. Découpage et classement : CoopMaths, Sébastien Lozano, Éric Elter et contributeurs. Adaptation de la présentation : Maths Explorer.
% CC BY-SA : https://coopmaths.fr/www/about/
% https://coopmaths.fr/alea/static/bac/2025/tex/bac_2025_06_sujet2_polynesie_4_cor.tex
\begin{document}
\noindent\textit{APMEP / CoopMaths. CC BY-SA. Présentation adaptée par Maths Explorer.}\par\medskip
\begin{enumerate}
	\item \textbf{Affirmation \theenumi : Fausse.}

	On a $\operatorname{Card}(E) = 7$ et $\operatorname{Card}(F) = 10$.

	Les 3-uplets d'éléments distincts de $E$, ce sont des arrangements de 3 éléments distincts (l'ordre compte, sans répétition) choisis parmi 7. Il y en a donc :

	$\dfrac{7 !}{(7 - 3)!} = \dfrac{7!}{4!} = 7 \times 6 \times 5 = 210$.

	Les combinaisons à 4 éléments de $F$ sont des ensembles (sans ordre, sans répétition) de 4 éléments choisis parmi 10. Il y en a :

	$\displaystyle \binom{10}{4} = \dfrac{10!}{(10-4)! \times 4!} = \dfrac{10 \times 9 \times 8 \times 7}{4 \times 3 \times 2} = 10 \times 3 \times 1 \times 7 = 210$.

	Il y a donc exactement le même nombre, et pas davantage, de 3-uplets d'éléments distincs de $E$ que de combinaisons à 4 éléments de $F$.

	\item \textbf{Affirmation \theenumi : Vraie.}

	En effet, le carré ABCD a un côté de 3 unités de longueurs, donc son aire est de :

	$3^2 = 9$ unités d'aire.

	Par ailleurs, la zone hachurée est délimitée par l'axe des abscisses, la courbe représentant la fonction carré et les droites d'équation $x = 0$ et $x=3$. La fonction carré étant positive, l'aire, en unités d'aire est - égale à l'intégrale :

$\displaystyle \int_{0}^{3} x^2 ~\mathrm{d} x = \left[\dfrac{1}{3}x^3\right]_0^3 = \left(\dfrac{1}{3}\times 3^3\right) - \left(\dfrac{1}{3}\times 0^3\right) = 9 - 0 = 9$ unités d'aire.

Les deux zones ont donc bien la même aire, de 9 unités d'aire.

	\item \textbf{Affirmation \theenumi : Fausse.}

Pour tout $x$ dans $[1~;~2]$, on pose :\quad
$\begin{cases}
u(x) = \dfrac{1}{2}x^2\\
v(x) = \ln(x)
\end{cases}$
\quad et \quad
$\begin{cases}
u'(x) = x\\
v'(x) = \dfrac{1}{x}
\end{cases}$

	On a alors : \quad
	$\begin{aligned}[t] \displaystyle
		J &= \int_{1}^{2} x\ln(x) ~\mathrm{dt} \\
		&= \int_{1}^{2} u'(x)\times v(x) ~\mathrm{dt} \quad \text{donc, par intégration par parties :}\\
		&= \Big[u(x) \times v(x)\Big]_1^2 - \int_{1}^{2} u(x)\times v'(x) ~\mathrm{dt} \\
		&= \left[\dfrac{1}{2}x^2 \times \ln(x) \right]_1^2 - \int_{1}^{2} \dfrac{1}{2}x^2\times \dfrac{1}{x} ~\mathrm{dt} \\
		&= \left(\dfrac{1}{2}\times 2^2 \times \ln(2)\right) - \left(\dfrac{1}{2}\times 1^2 \times \ln(1)\right) -\dfrac{1}{2} \int_{1}^{2} x~\mathrm{d}x\\
		&= 2 \ln(2) - 0 -\dfrac{1}{2} \left[\dfrac{1}{2}x^2\right]_{1}^{2}\\
		&= 2\ln(2) -\dfrac{1}{2}\left[\left(\dfrac{1}{2}\times 2^2\right) - \left(\dfrac{1}{2}\times 1^2\right)\right]\\
		&=2\ln(2) - \dfrac{3}{4} \approx \np{0,63629}
	\end{aligned}$

	Or $\dfrac{7}{11} \approx \np{0,63636}$, donc $J \neq \dfrac{7}{11}$.

	\item \textbf{Affirmation \theenumi : Vraie.}

D'une part, $f$ est dérivable sur $\R$ et, pour tout $x$ réel, on a :

$f'(x) = \e^{x} + 2\e^{2x}$.

D'autre part, pour tout réel $x$, on a :

$2f(x) - \e^{x} = 2\left(\e^{x} + \e^{2x}\right) - \e^{x} = 2\e^{x} + 2\e^{2x} - \e^{x} = \e^{x} + 2\e^{2x}$

On constate donc que, pour tout $x$ réel, on a : \quad $f'(x) = 2f(x) - \e^{x}$.

Autrement dit, la fonction $f$ vérifie l'équation différentielle $(E)$, donc $f$ est bien une solution de $(E)$.

	\item \textbf{Affirmation \theenumi : Vraie.}

Soit $x$ un réel dans $[0 ; 1[$. On a donc :\quad $(x - 1)\in [-1 ; 0[$.

Autrement dit, le réel $(x - 1)$ est \textbf{strictement négatif}.

Ainsi, comme on a :\quad $\lim\limits_{n \to +\infty} \e^{n} = +\infty$ \quad et que $(x-1)< 0$, on en déduit, par limite du produit que : \quad $\lim\limits_{n \to +\infty} (x-1) \e^{n} = -\infty$.

Par limite de la somme, on en déduit que : \quad $\lim\limits_{n \to +\infty} (x-1) \e^{n} + 1 = -\infty$.

Comme la fonction cos est bornée par 1 et $-1$, on a : \quad $\forall n \in \N,\quad -1 \leqslant \cos(n) \leqslant 1$,

et donc notamment :\quad $\forall n \in \N,\quad \cos(n) \leqslant 1 \implies (x-1) \e^{n} + \cos(n) \leqslant (x-1) \e^{n} + 1$.

Comme on a établi que :\quad $\lim\limits_{n \to +\infty} (x-1) \e^{n} + 1 = -\infty$ \quad cela implique, par comparaison, que :\quad $\lim\limits_{n \to +\infty} (x-1) \e^{n} + \cos(n) = -\infty$.

La suite $(u_n)$ diverge bien vers $-\infty$.
\end{enumerate}

\end{document}
